% !TEX root = ../main.tex

\section{Plug Detection and Following}
\label{sec:following}

\scaffold{
  Budget: about 180 words plus \cref{eq:position}.

  \textbf{Paragraph 1, plug detection.} A solve cannot tell an empty jack from a plugged cable with nothing behind it. The jacks have a \emph{tip switch}, a contact that touches the tip until a plug lifts it \cite{neutrik2026nmj6hfd2}. The plug check drives the tip low and the tip switch high: without a plug the pulls divide the rail (\qty{1.25}{\volt}), with one the tip switch reads its rail (\qty{2.5}{\volt}); above three quarters of the span counts as a plug. No current flows through the pedal, so the check works for any wiring and position.

  \textbf{Paragraph 2, classification.} After a solve, the arms decide what follows: a \emph{potentiometer} (exactly one arm within \qty{25}{\percent} of the total, the wiper), an \emph{end stop} or \emph{near end} (two small arms: the wiper rests on, or with a resistive lead near, a track end; the wiper is then taken from the per-jack setting or the wiper remembered for the jack), a \emph{tip-sleeve element} (switch or rheostat), nothing, or \emph{other}.

  \textbf{Paragraph 3, following.} An identified potentiometer is followed with one reading: the track ends are driven and only the wiper is read, which carries no current, so its contact resistance does not enter \cref{eq:position}. The track start is the sleeve whenever the sleeve is a track end, so the position rises away from the sleeve as in both common wirings. Every \qty{25}{\milli\second} one track end is re-read; if its drop changes or the wiper leaves the span, the jack is identified again. Switches and rheostats are followed through the tip (tip high, sleeve low); at most \qty{0.2}{\kilo\ohm} reads as closed, three readings between open and closed make a rheostat.
}

\begin{equation}
  p = \frac{V_\mathrm{w} - V_\mathrm{s}}{V_\mathrm{e} - V_\mathrm{s}}, \qquad 0 \leq p \leq 1
  \label{eq:position}
\end{equation}
